% !TEX root = ../matan.tex

\section{Квадратичная форма. Положительная и отрицательная определённость. Оценка снизу положительно определённой квадратичной формы. Достаточные условия экстремума}

\subsection*{Квадратичная форма и определённость}

\begin{Definition}{Квадратичная форма}{}
	\textbf{Квадратичной формой} на $\RR^d$ называется функция вида
	\[
	B(h)=\sum_{k=1}^{d}\sum_{j=1}^{d} a_{kj}\,h_k h_j,\qquad h=(h_1,\dots,h_d)\in\RR^d,
	\]
	где матрица $(a_{kj})$ симметрична ($a_{kj}=a_{jk}$). Форма однородна второй степени:
	$B(th)=t^2 B(h)$ для всех $t\in\RR$.
\end{Definition}

\begin{Definition}{Знакоопределённость}{}
	Квадратичная форма $B$ называется:
	\begin{itemize}
		\item \textbf{положительно определённой}, если $B(h)>0$ при всех $h\ne 0$;
		\item \textbf{отрицательно определённой}, если $B(h)<0$ при всех $h\ne 0$;
		\item \textbf{знакопеременной (неопределённой)}, если существуют $h_1,h_2$ с
		$B(h_1)>0$ и $B(h_2)<0$.
	\end{itemize}
\end{Definition}

\subsection*{Оценка снизу положительно определённой формы}

\begin{Lemma}{Оценка снизу}{}
	Если квадратичная форма $B$ положительно определена, то существует константа
	$\lambda>0$ такая, что
	\[
	B(h)\ge \lambda\,|h|^2\qquad\text{для всех } h\in\RR^d.
	\]
\end{Lemma}

\begin{proof}
	Квадратичная форма $B$ непрерывна на $\RR^d$ (это многочлен от координат).
	Рассмотрим единичную сферу $S=\{y\in\RR^d:|y|=1\}$. Сфера --- компакт (замкнута и
	ограничена), поэтому по теореме Вейерштрасса непрерывная функция $B$ достигает на
	ней наименьшего значения: существует $y_0\in S$ такое, что
	\[
	\lambda:=\min_{|y|=1}B(y)=B(y_0).
	\]
	Так как $y_0\ne 0$ и форма положительно определена, $\lambda=B(y_0)>0$.
	Для произвольного $h\ne 0$ вектор $h/|h|$ лежит на $S$, поэтому, пользуясь
	однородностью,
	\[
	B(h)=|h|^2\,B\!\left(\frac{h}{|h|}\right)\ge |h|^2\,\lambda=\lambda|h|^2.
	\]
	При $h=0$ неравенство тривиально. Лемма доказана.
\end{proof}

\begin{Remark}{Случай отрицательной определённости}{}
	Если $B$ отрицательно определена, то форма $-B$ положительно определена, и по
	лемме $-B(h)\ge\mu|h|^2$ с некоторым $\mu>0$, то есть $B(h)\le -\mu|h|^2$ для всех
	$h$.
\end{Remark}

\subsection*{Достаточные условия экстремума}

Пусть $f\colon\Omega\to\RR$ дважды дифференцируема в стационарной точке $x$
($\nabla f(x)=\vec 0$). Рассмотрим квадратичную форму второго дифференциала
\[
B(h)=\sum_{k=1}^{d}\sum_{j=1}^{d}\partial_k\partial_j f(x)\,h_k h_j,\qquad h\in\RR^d
\]
(её матрица --- матрица Гессе, симметрична по теореме Юнга).

\begin{Theorem}{Достаточное условие локального экстремума}{}
	В стационарной точке $x$:
	\begin{enumerate}
		\item если $B$ положительно определена, то $x$ --- точка строгого локального
		\textbf{минимума};
		\item если $B$ отрицательно определена, то $x$ --- точка строгого локального
		\textbf{максимума};
		\item если $B$ знакопеременна, то экстремума в точке $x$ \textbf{нет} (седловая
		точка).
	\end{enumerate}
\end{Theorem}

\begin{proof}
	Запишем формулу Тейлора второго порядка с остаточным членом в форме Пеано;
	обозначим $h=\xi-x$. Так как $\nabla f(x)=\vec 0$, линейная часть исчезает, и
	\[
	f(\xi)-f(x)=\tfrac12\sum_{k,j}\partial_k\partial_j f(x)\,h_k h_j+o(|h|^2)
	=\tfrac12 B(h)+o(|h|^2),\qquad h\to 0.
	\]
	Множитель $1/2$ на знак не влияет; для краткости включим его в $B$ (это не меняет
	знакоопределённости).

	\textbf{Случай 1: $B$ положительно определена.} По лемме $B(h)\ge\lambda|h|^2$ с
	$\lambda>0$. Тогда
	\[
	f(\xi)-f(x)=B(h)+o(|h|^2)\ge|h|^2\!\left(\lambda+\frac{o(|h|^2)}{|h|^2}\right).
	\]
	Дробь $o(|h|^2)/|h|^2\to 0$ при $h\to 0$, поэтому при достаточно малых $|h|\ne 0$
	скобка положительна, и $f(\xi)-f(x)>0$. Значит, $x$ --- строгий локальный минимум.

	\textbf{Случай 2: $B$ отрицательно определена.} Аналогично, $-B$ положительно
	определена, $B(h)\le-\mu|h|^2$ с $\mu>0$, и
	\[
	f(\xi)-f(x)\le|h|^2\!\left(-\mu+\frac{o(|h|^2)}{|h|^2}\right)<0
	\]
	при малых $|h|\ne 0$. Значит, $x$ --- строгий локальный максимум.

	\textbf{Случай 3: $B$ знакопеременна.} Возьмём единичные векторы $\zeta_1,\zeta_2$
	с $B(\zeta_1)>0$, $B(\zeta_2)<0$. Вдоль луча $\xi=x+t\zeta_1$:
	\[
	f(\xi)-f(x)=t^2 B(\zeta_1)+o(t^2)=t^2\!\left(B(\zeta_1)+\tfrac{o(t^2)}{t^2}\right)>0
	\]
	при малых $t\ne 0$; вдоль луча $\xi=x+t\zeta_2$ аналогично $f(\xi)-f(x)<0$.
	Таким образом, в любой окрестности $x$ есть точки как с бо\'{}льшим, так и с
	меньшим значением функции, --- экстремума нет.
\end{proof}

\begin{Remark}{Вырожденный случай}{}
	Если $B$ лишь \emph{неотрицательно} определена ($B(h)\ge 0$, но $B(h_0)=0$ для
	некоторого $h_0\ne 0$), достаточное условие не работает: знак приращения
	определяется членами высших порядков, и вопрос требует отдельного исследования.
\end{Remark}
